\section{Appendix E -- Journal-Core Bridge and Extraction Logic}
% R005_APPENDIX_READER_FRAME
Reader frame. This appendix supports the main manuscript by explaining corpus role, evidence class, audit route, and limitation boundary. It is not a detached raw dump.

\label{sec:oc13-journal-core-bridge}

This appendix explains how the journal-core article relates to the master monograph.

\subsection{Why the journal-core exists}

Some channels require a bounded, faster-reading scientific artifact.
Editors, reviewers, and external readers often need a paper-sized object whose defended claim can
be read in one sitting.
The journal-core exists for that purpose.
Its job is not to replace the monograph but to provide a narrow article that remains anchored to
the same scientific truth conditions.

\subsection{What the journal-core may omit}

The journal-core may compress:

\begin{itemize}
\item extensive chapter-scale didactic explanation,
\item large parts of the domain atlas,
\item broad module coverage,
\item detailed appendical inventories,
\item and some of the monograph-level pedagogical scaffolding.
\end{itemize}

It may not silently change:

\begin{itemize}
\item theorem scope,
\item explicit nonclaims,
\item claimed source basis,
\item figure meaning,
\item or the relation between inherited corpus and present defended claim.
\end{itemize}

\subsection{What the monograph provides to later artifacts}

Because the journal-core is extracted from the present volume, later owner-facing and public
channels do not need to improvise their own understanding of the science.
They can derive from the master monograph when they need depth, from the journal core when they
need bounded exposition, and from the companion layer when they need support evidence.

This is the practical publication value of the monograph.
It is not merely larger than the article.
It is the document that prevents every later projection from becoming a separate interpretive
exercise.

\clearpage
